\documentclass[10pt,a4paper]{amsart}
\usepackage{amsmath,amssymb,amsthm,graphicx}
% Retained minimal preamble for the Stacks Project chapter more-algebra.tex (More on Algebra).
% The chapter file pinned-original-source.tex carries only \input{preamble} on its first line, so no
% preamble text lives inside the chapter itself.  The Stacks Project uses one single shared preamble for
% every chapter; the snapshot stages it identically as sources/stacks-chow/preamble.tex and
% sources/stacks-etale-cohomology/preamble.tex.  The block below is copied mechanically, line for line,
% from sources/stacks-chow/preamble.tex lines 14-18 (the Xy-pic setup) and 154-217 (the native theorem
% declarations together with the mathematical macros).  Nothing was added, removed or redefined; the
% wrapper class and the amsmath/amssymb/amsthm/graphicx base packages are supplied by the assembler.
% Chapter: more-algebra.tex -- More on Algebra.
% For commutative diagrams we use Xy-pic
\usepackage[all]{xy}

% We use 2cell for 2-commutative diagrams.
\xyoption{2cell}

% Theorem environments.
%
\theoremstyle{plain}
\newtheorem{theorem}[subsection]{Theorem}
\newtheorem{proposition}[subsection]{Proposition}
\newtheorem{lemma}[subsection]{Lemma}

\theoremstyle{definition}
\newtheorem{definition}[subsection]{Definition}
\newtheorem{example}[subsection]{Example}
\newtheorem{exercise}[subsection]{Exercise}
\newtheorem{situation}[subsection]{Situation}

\theoremstyle{remark}
\newtheorem{remark}[subsection]{Remark}
\newtheorem{remarks}[subsection]{Remarks}

\numberwithin{equation}{subsection}

% Macros
%
\def\lim{\mathop{\mathrm{lim}}\nolimits}
\def\colim{\mathop{\mathrm{colim}}\nolimits}
\def\Spec{\mathop{\mathrm{Spec}}}
\def\Hom{\mathop{\mathrm{Hom}}\nolimits}
\def\Ext{\mathop{\mathrm{Ext}}\nolimits}
\def\SheafHom{\mathop{\mathcal{H}\!\mathit{om}}\nolimits}
\def\SheafExt{\mathop{\mathcal{E}\!\mathit{xt}}\nolimits}
\def\Sch{\mathit{Sch}}
\def\Mor{\mathop{\mathrm{Mor}}\nolimits}
\def\Ob{\mathop{\mathrm{Ob}}\nolimits}
\def\Sh{\mathop{\mathit{Sh}}\nolimits}
\def\NL{\mathop{N\!L}\nolimits}
\def\CH{\mathop{\mathrm{CH}}\nolimits}
\def\proetale{{pro\text{-}\acute{e}tale}}
\def\etale{{\acute{e}tale}}
\def\QCoh{\mathit{QCoh}}
\def\Ker{\mathop{\mathrm{Ker}}}
\def\Im{\mathop{\mathrm{Im}}}
\def\Coker{\mathop{\mathrm{Coker}}}
\def\Coim{\mathop{\mathrm{Coim}}}

% Boxtimes
%
\DeclareMathSymbol{\boxtimes}{\mathbin}{AMSa}{"02}

%
% Macros for moduli stacks/spaces
%
\def\QCohstack{\mathcal{QC}\!\mathit{oh}}
\def\Cohstack{\mathcal{C}\!\mathit{oh}}
\def\Spacesstack{\mathcal{S}\!\mathit{paces}}
\def\Quotfunctor{\mathrm{Quot}}
\def\Hilbfunctor{\mathrm{Hilb}}
\def\Curvesstack{\mathcal{C}\!\mathit{urves}}
\def\Polarizedstack{\mathcal{P}\!\mathit{olarized}}
\def\Complexesstack{\mathcal{C}\!\mathit{omplexes}}
% \Pic is the operator that assigns to X its picard group, usage \Pic(X)
% \Picardstack_{X/B} denotes the Picard stack of X over B
% \Picardfunctor_{X/B} denotes the Picard functor of X over B
\def\Pic{\mathop{\mathrm{Pic}}\nolimits}
\def\Picardstack{\mathcal{P}\!\mathit{ic}}
\def\Picardfunctor{\mathrm{Pic}}
\def\Deformationcategory{\mathcal{D}\!\mathit{ef}}
\usepackage{cleveref}
\crefname{originallemma}{original lemma}{original lemmas}
\Crefname{originallemma}{Original lemma}{Original lemmas}
\makeatletter
\newlabel{more-algebra-lemma-projective-amplitude}{{0A5P}{source}{}{sourcepointer-1}{}}
\newlabel{more-algebra-lemma-projective-amplitude@cref}{{[originallemma][0][]0A5P}{[1][0][]source}{}{}{}}
\makeatother
\crefname{originallemma}{original lemma}{original lemmas}
\Crefname{originallemma}{Original lemma}{Original lemmas}
\makeatletter
\newlabel{derived-lemma-subcategory-left-resolution}{{05T7}{source}{}{sourcepointer-2}{}}
\newlabel{derived-lemma-subcategory-left-resolution@cref}{{[originallemma][0][]05T7}{[1][0][]source}{}{}{}}
\makeatother
\crefname{originallemma}{original lemma}{original lemmas}
\Crefname{originallemma}{Original lemma}{Original lemmas}
\makeatletter
\newlabel{algebra-lemma-characterize-projective}{{05CF}{source}{}{sourcepointer-3}{}}
\newlabel{algebra-lemma-characterize-projective@cref}{{[originallemma][0][]05CF}{[1][0][]source}{}{}{}}
\makeatother
\title{A characterisation of bounded projective amplitude for a complex}
\author{The Stacks Project Authors}
\date{stacks-project:more-algebra.tex:local-snapshot-20261012; GNU FDL 1.2 or later}
\begin{document}
\maketitle
\setcounter{subsection}{1}
\noindent\textit{Editorial note.} All delivered mathematical text is the author's own, copied from the pinned Stacks Project chapter file more-algebra.tex; no statement or proof step was rewritten, reordered or removed.\par\medskip
\noindent\textit{Editorial note.} The chapter file pinned-original-source.tex carries only \textbackslash{}input\{preamble\} on its first line, so no preamble text lives inside the chapter. The retained file sources/stacks-more-algebra/retained-preamble.tex is a mechanical, line-for-line copy of the project's single shared preamble as staged in the snapshot (Xy-pic setup plus theorem declarations and mathematical macros). Nothing was added, removed or redefined, and the excerpt serves only as preamble, never as a span coordinate.\par\medskip
\noindent\textit{Editorial note.} Both counted spans start at the \textbackslash{}begin line of their environment, so the statement and the author's proof are delivered whole; surrounding prose of the section is deliberately outside the counted spans and is carried as standard prerequisites instead.\par\medskip
\noindent\textit{Editorial note.} The citation fields (source\_id, source\_version, authors, source\_url, license.name, source\_title) are identical to those of every other card this lane has produced from this chapter, so the batch selector treats them as one source.\par\medskip
\noindent\textit{Allowed standard prerequisites (editorial).}\begin{itemize}
\item Projective amplitude of a complex and the notion of a bounded representative. defined in the same section, source lines 17360-17412, before the counted statement. Supplies both sides of the characterisation.
\item Cohomology of a complex can be killed degree by degree by choosing a projective module mapping onto it. used at proof lines 17430-17460. Constructs the resolution one degree at a time.
\item Termination of the construction in the bounded case. used at proof lines 17460-17480. Shows that the resolution is bounded, which is what the characterisation requires.
\end{itemize}\medskip
\ifdefined\hypertarget\hypertarget{sourcepointer-1}{}\fi
\noindent\textit{Source reference.} Original lemma 0A5P: focused-seeds/AGP/sources/stacks-snapshot-20261011/tags.txt, source lines 12817--12817. The lemma's own canonical tag, quoted so that readers can resolve the label inside the delivered text; it points at the counted statement itself and carries no proof content.\par\medskip
\ifdefined\hypertarget\hypertarget{sourcepointer-2}{}\fi
\noindent\textit{Source reference.} Original lemma 05T7: focused-seeds/AGP/sources/stacks-snapshot-20261011/tags.txt, source lines 7471--7471. Navigation pointer to a statement the counted proof refers to but does not contain; nothing load-bearing is delegated to it.\par\medskip
\ifdefined\hypertarget\hypertarget{sourcepointer-3}{}\fi
\noindent\textit{Source reference.} Original lemma 05CF: focused-seeds/AGP/sources/stacks-snapshot-20261011/tags.txt, source lines 6918--6918. Navigation pointer to a statement the counted proof refers to but does not contain; nothing load-bearing is delegated to it.\par\medskip
\begin{lemma}
\label{lemma-projective-amplitude}
Let $R$ be a ring. Let $K$ be an object of $D(R)$. Let $a, b \in \mathbf{Z}$.
The following are equivalent
\begin{enumerate}
\item $K$ has projective-amplitude in $[a, b]$,
\item $\Ext^i_R(K, N) = 0$ for all $R$-modules $N$ and all
$i \not \in [-b, -a]$,
\item $H^n(K) = 0$ for $n > b$ and
$\Ext^i_R(K, N) = 0$ for all $R$-modules $N$ and all $i > -a$, and
\item $H^n(K) = 0$ for $n \not \in [a - 1, b]$ and
$\Ext^{-a + 1}_R(K, N) = 0$ for all $R$-modules $N$.
\end{enumerate}
\end{lemma}

\begin{proof}
Assume (1). We may assume $K$ is the complex
$$
\ldots \to 0 \to P^a \to P^{a + 1} \to \ldots \to
P^{b - 1} \to P^b \to 0 \to \ldots
$$
where $P^i$ is a projective $R$-module for all $i \in \mathbf{Z}$.
In this case we can compute the ext groups by the complex
$$
\ldots \to 0 \to \Hom_R(P^b, N) \to \ldots \to
\Hom_R(P^a, N) \to 0 \to \ldots
$$
and we obtain (2).

\medskip\noindent
Assume (2) holds. Choose an injection $H^n(K) \to I$ where $I$
is an injective $R$-module. Since $\Hom_R(-, I)$ is an exact functor,
we see that $\Ext^{-n}(K, I) = \Hom_R(H^n(K), I)$.
We conclude in particular that $H^n(K)$ is zero for $n > b$.
Thus (2) implies (3).

\medskip\noindent
By the same argument as in (2) implies (3) gives that (3) implies (4).

\medskip\noindent
Assume (4). The same argument as in (2) implies (3) shows that
$H^{a - 1}(K) = 0$, i.e., we have $H^i(K) = 0$ unless $i \in [a, b]$.
In particular, $K$ is bounded above and we can choose a
a complex $P^\bullet$ representing $K$ with $P^i$ projective
(for example free) for all $i \in \mathbf{Z}$ and $P^i = 0$ for $i > b$.
See Derived Categories, Lemma
\ref{derived-lemma-subcategory-left-resolution}.
Let $Q = \Coker(P^{a - 1} \to P^a)$. Then $K$ is quasi-isomorphic
to the complex
$$
\ldots \to 0 \to Q \to P^{a + 1} \to \ldots \to P^b \to 0 \to \ldots
$$
as $H^i(K) = 0$ for $i < a$.
Denote $K' = (P^{a + 1} \to \ldots \to P^b)$ the corresponding object of
$D(R)$. We obtain a distinguished triangle
$$
K' \to K \to Q[-a] \to K'[1]
$$
in $D(R)$. Thus for every $R$-module $N$ an exact sequence
$$
\Ext^{-a}(K', N) \to \Ext^1(Q, N) \to \Ext^{1 - a}(K, N)
$$
By assumption the term on the right vanishes. By the implication
(1) $\Rightarrow$ (2) the term on the left vanishes. Thus $Q$
is a projective $R$-module by
Algebra, Lemma \ref{algebra-lemma-characterize-projective}.
Hence (1) holds and the proof is complete.
\end{proof}

\end{document}
